\title{User Manual}
\def \documentid {SPW\_BRIDGE-UVIE-UM-001}
\date{Issue 1.1, September 10, 2026}

\newcommand\affil[1]{\textsuperscript#1}

\def\preparedby {Armin Luntzer\affil{1}}
\def\checkedby {Roland Ottensamer\affil{1}}
\def\approvedby {Franz Kerschbaum\affil{1}}

\def\affiliations{
	\affil{1} Department of Astrophysics, University of Vienna
}




\input{shared/template/univie.tex}
\input{shared/template/titlepage.tex}
\input{shared/template/traceability.tex}
\usepackage[nopostdot, numberedsection, style=super, section, toc, acronyms, nogroupskip]{glossaries}
\usepackage{xparse}

\setlength{\glsdescwidth}{0.8\textwidth}
\renewcommand{\glsnamefont}[1]{\textbf{#1}}

% label, acronym, name, description
\DeclareDocumentCommand{\newdualentry}{ m m m m } {
	\newglossaryentry{gls-#1}{
	name={#3 (\glstext{#1})},
	description={#4}, nonumberlist
	}
	\newglossaryentry{#1}{
		type=\acronymtype,
		name={#2},
		first={#3 (#2)},
		firstplural={#3s (#2s)},
		description=\glslink{gls-#1}{#3},
		nonumberlist
	}%
}
%%% \gls{ABBREV} to point at Acronym/Abbreviation index
%%% \gls{gls-ABBREV} to point to glossary

\makeglossaries

\newdualentry{ADC}%
  {ADC}%
  {Analog to Digital Converter}%
  {An Analog to Digital Converter is a system that converts an analog signal
   into a quantized digital signal. Its counterpart is the \gls{DAC}.}%

\newdualentry{ABI}%
  {ABI}%
  {Application Binary Interface}%
  {An Application Binary Interface defines the low-level interface between
   an application and the operating system or between application components
   running on the same hardware.}%

\newdualentry{API}%
  {API}%
  {Application Programming Interface}%
  {The Application Programming Interface defines how a developer can write
   a program that requests services from an operating system or application.
   \glspl{API} are implemented by function calls composed of verbs and nouns,
   i.e. a function to execute on an object.}%

\newdualentry{ASan}%
  {ASan}%
  {AddressSanitizer}%
  {AddressSanitizer is a dynamic analysis tool that detects memory errors in %
   C and C++ programs, such as buffer overflows, use-after-free, and memory %
   leaks, by instrumenting the compiler.}%

\newdualentry{BSP}%
  {BSP}%
  {Board Support Package}%
  {A Board Support Package is the implementation of a specific interface defined
   by the abstract layer of an operating system that enables the latter to run
   on the particular hardware platform.}%

\newdualentry{CERT}%
  {CERT}%
  {Computer Emergency Response Team}%
  {The Computer Emergency Response Team is the coordination centre at the %
   Software Engineering Institute that investigates and reports on software %
   security issues.}%

\newdualentry{CTU}%
  {CTU}%
  {Cross Translation Unit}%
  {Cross Translation Unit analysis is an extension to static analysis that %
   follows calls and references across translation unit boundaries rather than %
   analysing a single translation unit in isolation.}%

\newdualentry{CVSS}%
  {CVSS}%
  {Common Vulnerability Scoring System}%
  {The Common Vulnerability Scoring System is a standard for rating the %
   severity of software vulnerabilities.}%

\newdualentry{CWE}%
  {CWE}%
  {Common Weakness Enumeration}%
  {The Common Weakness Enumeration is a community developed list of common %
   software and hardware weakness types.}%

\newdualentry{CPU}%
  {CPU}%
  {Central Processing Unit}%
  {The Central Processing Unit is the electronic circuitry that interprets
  instructions of a computer program and performs control logic, arithmetic,
  and input/output operations specified by the instructions. It maintains
  high-level control of peripheral components, such as memory and other devices.}%

\newdualentry{DAC}%
  {DAC}%
  {Digital to Analog Converter}%
  {A Digital to Analog Converter is a system that converts a quantized digital
   signal into an analog signal. Its counterpart is the \gls{ADC}.}%

\newdualentry{DMA}%
  {DMA}%
  {Direct Memory Access}%
  {Direct Memory Access is a feature of a computer system that allows hardware
   subsystems to access main system \gls{gls-RAM} directly, thereby bypassing
   the \gls{gls-CPU}.}%

\newacronym{DRD}{DRD}{Document Requirements Definition}

\newacronym{CIDL}{CIDL}{Configuration Item Data List}

\newacronym{ADD}{ADD}{Architectural Design Document}

\newacronym{CCB}{CCB}{Configuration Control Board}

\newacronym{IRD}{IRD}{Interface Requirements Document}

\newacronym{PQA}{PQA}{Product Quality Assurance Plan}

\newacronym{QTS}{QTS}{Qualification Test Specification}

\newdualentry{DSP}%
  {DSP}%
  {Digital Signal Processor}%
  {A Digital Signal Processor is a specialised processor with its architecture %
   targeting the operational needs of digital signal processing.}%

\newglossaryentry{Doxygen}{
  name={Doxygen},
  description={Doxygen is a documentation generator for writing software %
   	       reference documentation.},
  nonumberlist
  }

\newdualentry{ELF}%
  {ELF}%
  {Executable and Linkable Format}%
  {The Executable and Linkable Format is a common standard file format for
   executables, object code, shared libraries, and core dumps.}%

\newacronym{FDIR}{FDIR}{Fault Detection, Isolation and Recovery}

\newdualentry{EDAC}%
  {EDAC}%
  {Error Detection And Correction}%
  {Error Detection And Correction is a technique used to detect and correct
   errors in digital data, typically in memory, through the use of redundant
   check bits. On fault-tolerant SPARC processors it is used to protect memory
   against single-event upsets.}%

\newdualentry{EDF}%
  {EDF}%
  {Earliest Deadline First}%
  {Earliest Deadline First is a dynamic priority scheduling algorithm that
   assigns priorities based on task deadlines. Tasks with the nearest
   deadline receive the highest priority.}%

\newdualentry{FIFO}%
  {FIFO}%
  {First In - First Out}%
  {In FIFO processing, the "head" element of a queue is processed first.
   Once complete, the element is removed and the next element in line becomes
   the new queue head.}%

\newglossaryentry{FlightOS}{
  name={FlightOS},
  description={FlightOS is a lightweight real-time operating system for
               space-based astronomy applications, targeting European
               \gls{SPARC}-based processors. It is licensed under the
               \gls{GPL}.},
  nonumberlist
  }

\newdualentry{FPU}%
  {FPU}%
  {Floating Point Unit}%
  {A co-processor unit that specialises in floating-point calculations.}%

\newdualentry{ILP}%
  {ILP}%
  {Instruction Level Parallelism}%
  {Instruction-level parallelism (ILP) is a measure of how many instructions in
   a computer program can be executed simultaneously by the \gls{CPU}.}%

\newdualentry{ISR}%
  {ISR}%
  {Interrupt Service Routine}%
  {An Interrupt Service Routine is a function that handles the actions needed
   to service an interrupt.}%

\newacronym{FSW}{FSW}{Flight Software}

\newdualentry{GCC}%
  {GCC}%
  {GNU Compiler Collection}%
  {The GNU Compiler Collection is a compiler system produced by the
   GNU project. It is part of the GNU toolchain collection of programming
   tools.}%

\newdualentry{GPL}%
  {GPL}%
  {GNU General Public License}%
  {The GNU General Public License is a free software licence. Version 2 %
   (GPLv2) is used to license FlightOS and, unless stated otherwise, the %
   documentation examples reproduced in this document.}%

\newglossaryentry{GNU}{
  name={GNU},
  description={GNU is a recursive acronym that stands for "GNU's Not Unix!".
	      The GNU project is a free software collaboration project announced
	      in 1987. Users are free to run GNU-licensed software, share, copy,
	      distribute, study and modify it. GNU software guarantees these
	      freedom-rights legally via its license and is thefore free
      	      software.},
  nonumberlist
}

\newglossaryentry{LEON2}{
  name={LEON2},
  description={The LEON2 is a synthesisable VHDL model of a 32-bit processor
	       compliant with the SPARC V8 architecture. It is highly
	       configurable and particularly suitable for \gls{SoC} designs.
	       Its source code is available under the GNU LGPL license},
  nonumberlist
  }

\newglossaryentry{LEON3}{
  name={LEON3},
  description={The LEON3 is an updated version of the \gls{LEON2}, changes
	       include \gls{gls-SMP} support and a deeper instruction pipeline},
  nonumberlist
  }

\newglossaryentry{LEON3-FT}{
  name={LEON3-FT},
  description={The LEON3-FT is a fault-tolerant version of the \gls{LEON3}.
  	       Changes to the base version include autonomous error handling,
       	       cache locking and different cache replacement strategies.},
  nonumberlist
  }

\newglossaryentry{LEON4}{
  name={LEON4},
  description={The LEON4 is an updated version of the \gls{LEON3} family of
               SPARC V8 processors, offering a deeper pipeline, higher clock
               frequencies and a wider memory interface.},
  nonumberlist
  }

\newglossaryentry{GR740}{
  name={GR740},
  description={The GR740 is a system-on-chip based on four \gls{LEON4FT}
               cores with \gls{SMP} support, manufactured by Gaisler
               \cite{GR740UM}.},
  nonumberlist
  }

\newglossaryentry{GR712RC}{
  name={GR712RC},
  description={The GR712RC is a dual-core \gls{LEON3-FT} system-on-chip
               manufactured by Gaisler \cite{GR712UM}.},
  nonumberlist
  }

\newglossaryentry{LEON4FT}{
  name={LEON4FT},
  description={The LEON4FT is the fault-tolerant version of the \gls{LEON4}
               processor core, adding error detection and correction
               capabilities.},
  nonumberlist
  }

\newglossaryentry{ramses}{
  name={RAMSES},
  description={The RAMSES project is a space mission whose on-board data
               processing platform is based on a custom built computing device
               using the \gls{GR712RC} processor.},
  nonumberlist
  }

\newglossaryentry{ariel}{
  name={ARIEL},
  description={The ARIEL (Atmospheric Remote-sensing Infrared Exoplanet
               Large-survey) project is a space mission whose on-board data
               processing platform is based on a custom built computing device
               using the \gls{GR712RC} processor.},
  nonumberlist
  }

\newglossaryentry{smile}{
  name={SMILE},
  description={The SMILE (Solar wind Magnetosphere Ionosphere Link Explorer)
               project is a space mission whose on-board data processing
               platform is based on a custom built computing device using the
               \gls{GR712RC} processor.},
  nonumberlist
  }

\newdualentry{MMU}%
  {MMU}%
  {Memory Management Unit}%
  {A Memory Management Unit performs address space translation between physical
   and virtual memory pages and protects unprivileged access to certain memory
   regions.}%

\newdualentry{MPPB}%
  {MPPB}%
  {Massively Parallel Processor Breadboarding system}%
  {The Massively Parallel Processor Breadboarding system is a proof-of-concept %
   design for a space-hardened, fault-tolerant multi-DSP system with various %
   subsystems to build a powerful digital signal processing system with a high %
   data throughput. Its distinguishing features are the \gls{gls-NoC} and the
   \gls{Xentium} \glspl{DSP} controlled by a \gls{LEON2} processor.
   It was developed under ESA contract 21986 by Recore Systems B.V.}%

\newacronym{NCR}{NCR}{Non-Conformance Reports}

\newdualentry{NGAPP}%
  {NGAPP}%
  {Next Generation Astronomy Processing Platform}%
  {Next Generation Astronomy Processing Platform was an evaluation of the
   \gls{MPPB} performed in a joint effort of RUAG Space Austria and the
   Department of Astrophysics of the University of Vienna.
   The project was funded under ESA contract 40000107815/13/NL/EL/f.}%

\newdualentry{NoC}%
  {NoC}%
  {Network On Chip}%
  {A Network On Chip is a communication system on an integrated circuit that
   applies (packet based) networking to on-chip communication. It offers
   improvements over more conventional bus interconnects and is more scalable
   and power efficient in complex \gls{gls-SoC} desgins.}%

\newacronym{NRB}{NRB}{Nonconformance Review Board}

\newacronym{PA}{PA}{Product Assurance}

\newacronym{PAM}{PAM}{Product Assurance Manager}

\newdualentry{POSIX}
  {POSIX}
  {Portable Operating System Interface}
  {The Portable Operating System Interface is a family of standards specified
   by the IEEE Computer Society for maintaining compatibility between
   operating systems.}%

\newdualentry{PUS}
  {PUS}
  {Packet Utilisation Standard}
  {The Packet Utilisation Standard addresses the end-to-end transport of telemetry
   and telecommand data between user applications on the ground and applications
   onboard a satellite. See also ECSS-E-70-41A.}%

\newdualentry{RAM}%
  {RAM}%
  {Random-Access Memory}%
  {Random-Access Memory is a type of memory where each memory cell may be
   accessed directly via their memory addresses.}%

\newacronym{RID}{RID}{Review Item Discrepancy}

\newdualentry{RISC}%
  {RISC}%
  {Reduced Instruction Set Computing}%
  {RISC is a \gls{CPU} design strategy that intends to improve performance by
   combining a simplified instruction set with a microprocessor architecture
   that is capable of executing an instruction in a smaller number of clock
   cycles.}%

\newdualentry{RMAP}%
  {RMAP}%
  {Remote Memory Access Protocol}%
  {The Remote Memory Access Protocol is a form of \gls{SpaceWire} communication
   that transparently communicates writes to memory mapped regions between
   different hardware devices.}%


\newdualentry{ROSE}%
  {ROSE}%
  {ROSE compiler framework}%
  {The ROSE compiler framework is an open source compiler infrastructure that %
   provides a rich abstract syntax tree representation of source code, on which %
   static analysis tools such as Rosecheckers are built.}%

\newdualentry{RR}%
  {RR}%
  {Round Robin}%
  {Round Robin is a scheduling algorithm where time slices are assigned in equal
   poritions and in circular order. In the context of threads, priorities are
   usually only used to control re-scheduling order when a mutex is accessed by
   a thread.}%

\newdualentry{RTOS}%
  {RTOS}%
  {Real-Time Operating System}%
  {A Real-Time Operating System is an operating system intended to serve
   real-time applications that process data as it comes in, typically
   without buffer delays, and guarantees a certain capability within a
   specified time constraint.}%

\newacronym{RSA}{RSA}{RUAG Space Austria}

\newacronym{SCF}{SCF}{Software Configuration File}

\newacronym{SDD}{SDD}{Software Design Document}

\newacronym{SDP}{SDP}{Software Development Plan}

\newacronym{SMCP}{SMCP}{Software Management and Configuration Plan}

\newacronym{SPSC}{SPSC}{Single-Producer Single-Consumer}

\newdualentry{SMP}%
  {SMP}%
  {Symmetric Multiprocessing}%
  {Symmetric Multiprocessing denotes computer architectures, where two or more
   identical processors are connected to the same periphery and are controlled
   by the same operating system instance.}%

\newdualentry{SEI}%
  {SEI}%
  {Software Engineering Institute}%
  {The Software Engineering Institute is a research and development centre at %
   Carnegie Mellon University that advances software engineering practice and %
   publishes the SEI CERT Coding Standards.}%

\newacronym{SPAMR}{SPAMR}{Software Product Assurance Milestone Report}

\newacronym{SPR}{SPR}{Software Problems Reports}

\newacronym{SSS}{SSS}{System Software Specification}

\newacronym{SRD}{SRD}{Software Release Document}

\newacronym{SRS}{SRS}{Software Requirement Specification}

\newdualentry{SoC}%
  {SoC}%
  {System On Chip}%
  {A System On Chip is an integrated circuit that combines all components of a %
   computer or other electronic system into a single chip.}%

\newdualentry{SRMMU}%
  {SRMMU}%
  {SPARC Reference Memory Management Unit}%
  {The SPARC Reference Memory Management Unit is the reference implementation
   of an MMU for the \gls{SPARC} architecture. It provides virtual memory
   translation and memory protection.}%

\newglossaryentry{SpaceWire}{
  name={SpaceWire},
  description={SpaceWire is a spacecraft communication network based in part
               on the IEEE 1355 standard of communications.},
  nonumberlist
  }

\newglossaryentry{SPARC}{
  name={SPARC},
  description={SPARC ("scalable processor architecture") is a \gls{gls-RISC}
  	       instruction set architecture developed by Sun Microsystems in
	       the 1980s. The distinct feature of SPARC processors is the high
	       number of \gls{gls-CPU} registers that are accessed similarly to
	       stack variables via ``sliding windows''.},
  nonumberlist
  }


\newacronym{SQ}{SQ}{Software Quality}

\newacronym{SQAM}{SQAM}{Software Quality Assurance Manager}

\newdualentry{SSDP} % label
  {SSDP}            % abbreviation
  {Scalable Sensor Data Processor}  % long form
  {The Scalable Sensor Data Processor (SSDP) was a next generation on-board %
   data processing mixed-signal ASIC, envisaged for use in scientific payloads %
   requiring high-performance on-board processing capabilities. It was built on %
   a heterogeneous multicore architecture, combining two %
   \gls{Xentium} \gls{DSP} cores with a general-purpose \gls{LEON3-FT} control %
   processor in a \gls{gls-NoC}.} % description

\newacronym{SQA}{SQA}{Software Quality Assurance}

\newdualentry{TCM}%
  {TCM}%
  {Tightly-Coupled Memory}%
  {Tightly-Coupled Memory is the local data memory that is directly accessible %
   by a Xentium's load/store unit. It can be viewed as a completely %
   program-controlled data cache.}%

\newdualentry{TOCTOU}%
  {TOCTOU}%
  {Time Of Check To Time Of Use}%
  {Time Of Check To Time Of Use is a vulnerability where a resource is checked
   for a condition at one point in time and then used at a later point in time,
   when the condition may have changed.}%


\newacronym{TP}{TP}{Test Plan}
\newacronym{TS}{TS}{Test Specification}

\newacronym{SVR}{SVR}{Software Verification Report}

\newacronym{VQR}{VQR}{Software Verification and Quality Report}

\newdualentry{UBSan}%
  {UBSan}%
  {UndefinedBehaviorSanitizer}%
  {UndefinedBehaviorSanitizer is a fast undefined behavior detector that %
   modifies the program at compile-time to catch various kinds of undefined %
   behavior during program execution.}%

\newdualentry{UVIE}%
  {UVIE}%
  {University of Vienna}%
  {University of Vienna.}%

\newdualentry{VLIW}%
  {VLIW}%
  {Very Long Instruction Word}%
  {Very Long Instruction Word is a processor architecture design concept that
   exploits \gls{gls-ILP}. This approach allows higher performance at a smaller
   silicone footprint compared to serialised instruction processors, as no
   instruction re-ordering logic to exploit superscalar capabilities of the
   processor must be integrated on the chip, but requires either code to be
   tuned manually or a very sophisticated compiler to exploit the full potential
   of the processor.}%

\newdualentry{WCET}%
  {WCET}%
  {Worst-Case Execution Time}%
  {Worst-Case Execution Time is the maximum length of time that a task or
   function can take to execute on a specific hardware platform. It is a
   critical parameter in the design and analysis of real-time systems.}%

\newglossaryentry{Xentium}{
  name={Xentium},
  description={The Xentium is a high performance \gls{gls-VLIW} \gls{DSP} core.
               It operates 10 parallel execution slots supporting 32/40 bit
	       scalar and two 16-bit element vector operations.},
  nonumberlist
  }


\newdualentry{CCSDS}%
  {CCSDS}%
  {Consultative Committee for Space Data Systems}%
  {The Consultative Committee for Space Data Systems is the body that publishes %
   the Space Packet Protocol used as the base transport of \gls{PUS} packets.}%

\newdualentry{CRC}%
  {CRC}%
  {Cyclic Redundancy Check}%
  {A Cyclic Redundancy Check is an error-detecting code commonly used to detect %
   accidental changes to digital data, e.g. in the \gls{RMAP} and \gls{PUS} %
   protocols.}%

\newdualentry{EGSE}%
  {EGSE}%
  {Electrical Ground Support Equipment}%
  {Electrical Ground Support Equipment is the collection of test and support %
   equipment used to verify a space segment on ground, e.g. the CHEOPS model %
   operated through the SpaceWire bridge.}%

\newdualentry{EOP}%
  {EOP}%
  {End Of Packet}%
  {End Of Packet is the \gls{SpaceWire} marker that terminates a packet, the %
   counterpart of the Error End Of Packet (EEP) marker that signals an %
   aborted transmission.}%

\newdualentry{FEE}%
  {FEE}%
  {Front-End Electronics}%
  {Front-End Electronics is the electronics directly interfacing detectors and %
   processing their signals; on CHEOPS the FEE produces the housekeeping and %
   data frames that the bridge can interpret.}%

\newglossaryentry{GRESB}{
  name={GRESB},
  description={The GRESB is a SpaceWire/Ethernet bridge by Gaisler that routes %
	       \gls{SpaceWire} packets to and from TCP sockets over an Ethernet %
	       network; it offers three SpaceWire links and six virtual links %
	       \cite{GRESBUM}.},
  nonumberlist
  }

\newglossaryentry{STAR}{
  name={STAR API},
  description={The STAR API is the software interface to SpaceWire interface %
	       devices such as the Brick MK II, Brick MK III, PCI and PCIe %
	       cards; \texttt{spw\_bridge} uses it for all device, channel and %
	       link management.},
  nonumberlist
  }

\newglossaryentry{SpW}{
  name={SpW},
  description={SpW is a common abbreviation of \gls{SpaceWire} used throughout %
	       this document.},
  nonumberlist
  }

\newdualentry{TCP}%
  {TCP}%
  {Transmission Control Protocol}%
  {Transmission Control Protocol is the stream-oriented transport protocol of %
   the Internet protocol suite; the bridge carries \gls{SpaceWire} traffic %
   over a single TCP connection by default.}%

\newdualentry{UDP}%
  {UDP}%
  {User Datagram Protocol}%
  {User Datagram Protocol is the datagram-oriented transport protocol of the %
   Internet protocol suite; with \texttt{-u} the bridge maps a single UDP %
   datagram to a single \gls{SpaceWire} packet.}%

\newdualentry{ECSS}%
  {ECSS}%
  {European Cooperation for Space Standardization}%
  {The European Cooperation for Space Standardization issues the standards %
   series used by the European space industry, e.g. ECSS-E-ST-70-41C for the %
   \gls{PUS} and ECSS-E-ST-50-52C for \gls{RMAP}.}%

\newglossaryentry{MSSL}{
  name={MSSL},
  description={The Mullard Space Science Laboratory is a department of %
	       University College London; on CHEOPS it provides the science %
	       payload, whose FEE data frames follow MSSL-IF-109 %
	       \cite{MSSLIF109}.},
  nonumberlist
  }


% do not remove
\glsresetall
\makeglossaries


\usepackage{vhistory}

\usepackage[giveninits=true]{biblatex}
\addbibresource{shared/bibliography.bib}
\nocite{ECSS5012C,ECSS5052C,ECSS7041C,ECSS7041A,CCSDS133,GRESBUM,GRMON2UM,MSSLIF109,GR712UM,GR740UM}

\usepackage{listings}

\lstdefinestyle{customc}{
  belowcaptionskip=1\baselineskip,
  breaklines=true,
  xleftmargin=\parindent,
  language=C,
  showstringspaces=false,
  basicstyle=\footnotesize\ttfamily,
  keywordstyle=\bfseries\color{uvie-orangered!75},
  commentstyle=\itshape\color{uvie-gray},
  identifierstyle=\color{uvie-blue},
  stringstyle=\color{uvie-goldenyellow},
}

\usepackage{appendix}

\begin{document}
\lstset{style=customc} 


\setmainfont{MyriadPro-SemiCondensed}
\uvietitlepage%
{SPW\_bridge --\\ A SpaceWire to network bridge for ground support equipment}%
{\doctitle}%
{shared/images/logo2.pdf}%
{}\setmainfont{MyriadPro}

\approvalpage

\tableofcontents
\newpage



\begin{versionhistory}
  \vhEntry{1.0}{10.09.2026}{AL}{Initial version}
  \vhEntry{1.1}{10.09.2026}{AL}{add monitor mode bridging two SpaceWire channels with RMAP decoding, payload suppression and a short debug form}
\end{versionhistory}


\chapter{Introduction}

\section{Purpose of the Document}

This document gives an overview on how to build, configure and use the
\gls{SpW} bridge software. An overview of the structure of the software will be
given; detailed information relevant to system programming must however be
looked up in the in-line documentation of the source files. The main text
describes the generic bridging machinery that the software provides, the
command-line interface, and the packet formats exchanged with a peer on the
network side.

\chapter{Applicable and Reference Documents}
\label{cha:appdoc}

\printbibliography[heading=none]


\chapter{Terms, Definitions and Abbreviated Items}
\printglossary[type=acronym]
\printglossary[type=main, style=altlist]



\chapter{Software Overview}

\section{Function and Purpose}

\texttt{spw\_bridge} is a software bridge between a \gls{SpaceWire} link and a
network connection to a workstation. It operates a \gls{SpaceWire} interface
device through the \gls{STAR} API and relays packets between the \gls{SpW}
link and a \gls{TCP}/IP or \gls{UDP} endpoint, in either direction.
\\

On the \gls{SpW} side, the software uses the STAR API to select a device and
channel, to set the link signalling rate and to submit and receive packets via
the STAR transfer-operation machinery. On the network side, it provides a mult
iple-connection server, a single-connection client, or a UDP endpoint. The
bridge may optionally interpret the carried data stream: \gls{PUS} packets,
\gls{FEE} data frames, GRESB protocol packets, and a pseudo-\gls{RMAP} command
interface are supported and are described below.
\\

The intended use is ground support equipment, e.g. to operate the front-end
electronics of an instrument through a \gls{SpaceWire}-to-Ethernet path from a
workstation, or to debug a \gls{LEON3-FT}-based target such as a
\gls{GR712RC} through the \gls{RMAP} protocol.

\section{Build Requirements}

In order to build \texttt{spw\_bridge}, make sure you have set up your
environment so that the STAR API tools are available. The build currently
links against the shared libraries

\begin{lstlisting}[language=bash]
  -lstar-api -lrmap_packet_library -lstar_conf_api_generic
  -lstar_conf_api_mk2 -lstar_conf_api_brick_mk2
  -lstar_conf_api_brick_mk3 -lstar_conf_api_pcie_mk2
\end{lstlisting}

\noindent
together with \texttt{-lpthread}. The matching headers are found under
\texttt{star-system/inc/star} and are provided by the STAR API installation.
You will also need a \gls{GCC} compiler for your host platform, GNU make and a
bash shell, i.e. what you would typically find in a Linux distribution with a
basic set of development tools installed. If you want to generate source code
documentation, you will also need to install \gls{Doxygen}.

\noindent
The \gls{STAR} API, the RMAP packet library and the device configuration
libraries are supplied by STAR-Dundee Ltd.\ as part of the support software of
their \gls{SpaceWire} devices; obtain them from STAR-Dundee if you are a
customer who acquired one of their devices.

\section{Build System}

The build is driven by the plain Makefile in the source root. To build the
binary, simply issue

\begin{lstlisting}[language=bash]
  $ make
\end{lstlisting}

\noindent
on your shell prompt from the top-level \texttt{spw\_bridge} directory. The
build produces an executable named \emph{spw\_bridge} in the current
directory. To remove the build output, use

\begin{lstlisting}[language=bash]
  $ make clean
\end{lstlisting}

\noindent
The optional \texttt{DEBUG} variable controls the verbosity of the debug
output. By default it is set to 1, which disables the per-packet \texttt{DBG}
traces. Setting \texttt{DEBUG} to 2 or higher enables them:

\begin{lstlisting}[language=bash]
  $ make DEBUG=2
\end{lstlisting}

\section{Source and Support Documentation}

The source files carry Doxygen-style in-line documentation. The \gls{STAR} API
header set and its reference documentation are kept below
\texttt{star-system/}, with the generated \gls{Doxygen} HTML under
\texttt{star-system/doc}. The reference standards drawn upon by the packet
handling code are listed in Chapter~\ref{cha:appdoc}.

\section{Current Platform Configuration}

The bridge is a host-side application and runs on a Linux workstation. The
\gls{STAR} device selected by default is device number 0; the device
enumeration performed at startup lists the detected devices and their names.
A channel number between 0 and 31 may be selected; the default channel is 1.
The link signalling rate is set to 10\,Mbit/s by default and may be configured
between 2 and 400\,Mbit/s. Channels, link identifiers and routing paths are
under configuration control through the command-line options described in
Chapter~\ref{cha:usage}.

\section{Features Overview}

The following gives a concise overview of the currently implemented features:

\begin{itemize}

\item Network transport over a single TCP connection (server or client mode)
or over UDP datagrams.\\

\item Interpretation of the network byte stream into \gls{PUS} packets
(ECSS-E-ST-70-41C) with sequence counting and CRC16 verification in both
directions.\\

\item Decoded printing of \gls{CCSDS} and PUS-C header fields, the per-packet
CRC16 status and the payloads for debugging.\\

\item Interpretation of FEE data frames (with the FEE framing of
\gls{MSSL} Draft A.14).\\

\item Support for the GRESB protocol, intended for use with the grmon debug
monitor for \gls{RMAP} access to a target.\\

\item A pseudo-\gls{RMAP} command service on a dedicated TCP port for
lossless read and write access to the \gls{SpW} network.\\

\item Optional stripping of a fixed number of leading bytes of incoming
\gls{SpW} packets (e.g. a path header added by the network).\\

\item Optional throttling of \gls{SpW} packet submission by an inserted delay
between packets.\\

\item Optional device reset on startup.\\

\item Monitor mode: verbatim copying of packets between two \gls{SpW} channels
of the same device, with the network side acting as a passive observer.\\

\end{itemize}


\chapter{Structure and Concept}

\section{Source Code Organisation}

The bridge consists of a small set of dedicated source files:

\begin{tabu} to \linewidth {@{}ll@{}}
	\texttt{spw\_bridge.c}	& main program: option processing, startup and shutdown \\
	\texttt{spw\_bridge.h}	& shared configuration and subsystem interfaces \\
	\texttt{net.c}		& network subsystem: sockets, framing, packet sink \\
	\texttt{spw.c}		& SpaceWire subsystem: STAR API device, link handling and the monitor mode packet sink \\
	\texttt{kbd.c}, \texttt{kbd.h}& keyboard monitor: toggles the debug printout and the short debug form \\
	\texttt{rmap.c}, \texttt{rmap.h}& RMAP command and reply helper functions \\
	\texttt{gresb.c}, \texttt{gresb.h}& GRESB protocol packet helpers \\
	\texttt{byteorder.h}	& endianness conversion macros \\
	\texttt{debug.h}		& debug output control \\
	\texttt{ecss/}		& \gls{CCSDS}, \gls{ECSS} (\gls{PUS}) and \gls{CRC} packet library \\
\end{tabu}

\noindent
The two subsystems talk to each other through the \texttt{bridge\_cfg}
structure: the network code installs a packet sink callback
(\texttt{cfg->pkt\_sink = net\_pkt\_sink}) that the \gls{SpW} receive loop
invokes for every complete packet received on the link.

\section{Network Subsystem}
\label{sec:net}

The network subsystem maintains a per-service fds set and multiplexes its
connections with \texttt{select()}. Two services exist: the \emph{data}
service carries the forwarded \gls{SpW} packets, and the \emph{rmap} service
implements the pseudo-\gls{RMAP} interface (see Section~\ref{sec:rmapif}).
Each service has its own accept and poll thread; the RMAP service is always
started as a listening server, even when the data path runs in client mode.
A peer that stays silent for more than 30 seconds is dropped.
\\

Depending on the selected mode, a received network byte stream is framed as
follows:

\begin{itemize}

\item \emph{raw}: one \texttt{recv()} call defines one \gls{SpW} packet. Note
that back-to-back packets sent on a stream socket may be merged into a single
receive, and a partially delivered packet may be split, so a higher-level
framing protocol is normally required for reliable transport.\\

\item \emph{PUS} (\texttt{-P}): a six-byte CCSDS primary header is read
first; the packet length field plus one octet plus the six-byte header
defines the packet boundary. The sequence counter in bytes 2 and 3 is checked
against the running sequence in both directions and mismatches are reported
on {\ttfamily stderr}.\\

\item \emph{FEE} (\texttt{-F}): a ten-byte FEE data header is expected; if
its protocol identifier is \texttt{0xF0}, the big-endian data length field
plus the header defines the packet boundary. A different protocol identifier
defers to raw framing.\\

\item \emph{GRESB} (\texttt{-G}): a four-byte host-to-GRESB header (protocol
byte plus three-byte big-endian size) defines the packet boundary.\\

\end{itemize}

In the other direction, every packet received on the \gls{SpW} link is passed
to \texttt{net\_pkt\_sink}, which forwards it to all connected data clients.
In GRESB mode the \gls{SpW} packet is wrapped into a host-to-GRESB data packet
before transmission.

\section{SpaceWire Subsystem}

The \gls{SpW} subsystem selects the STAR device and channel, optionally resets
the device, sets the link signalling rate through
\texttt{CFG\_setTransmitSignallingRate}, opens the channel and creates the
routing path with \texttt{STAR\_createAddress}. The link is then forced into
the running state and the measured receive rate is reported.
\\

A dedicated polling thread is started which repeatedly submits a single receive
transfer operation (\texttt{STAR\_createRxOperation} with \texttt{EOP}
terminated packets) and waits for its completion. On completion the received
packet is handed to the network sink. Transmission uses
\texttt{STAR\_createPacket} plus a transmit transfer operation, with an
optional throttle delay after each packet. Reads and writes of registers on
remote targets are implemented with the STAR RMAP packet library, see below.
\\
In monitor mode (Section~\ref{sec:monitor}) two channels of the device are
opened instead of one, each polled by its own thread; received packets are
copied to the other channel without a routing header.

\section{RMAP Helper Functions}

The RMAP helper sources provide verification, CRC-8 computation, and parsing
of RMAP packets, including a table of the command codes and status codes of
\gls{ECSS} ECSS-E-ST-50-52C. The bridge uses these helpers to parse RMAP
packets for logging and to decode replies received from targets on the
\gls{SpW} link.

\section{GRESB Support}

The GRESB is a \gls{SpaceWire}-to-Ethernet bridge intended for ground
equipment (see the \gls{GRESB} glossary entry and
\cite{GRESBUM}). The bridge implements its host protocol: a data transfer
consists of a one-byte protocol identifier followed by a three-byte big-endian
size field and the raw \gls{SpW} data. The reverse direction carries a header
with reserved, truncated and \gls{EOP}-error bits followed by the three-byte
size and the data.
\\

The GRESB uses a base port number of 3000 with two ports per virtual link:
the even port is the transmit port and the odd port is the receive port. Up
to five virtual links are supported. This mode is intended to be used with
the grmon debug monitor for \gls{RMAP} access (see
Section~\ref{sec:gresb}).

\section{ECSS Packet Library}

The \texttt{ecss/} directory contains a self-contained packet library that
follows the CCSDS Space Packet Protocol \cite{CCSDS133} and the \gls{PUS}
\cite{ECSS7041C}:

\begin{itemize}

\item \texttt{ccsds\_pkt.c/h} -- field accessors and setters for the
six-byte CCSDS primary header (version, type, secondary header flag, APID,
sequence flags and count, data length).\\

\item \texttt{ecss\_pkt.c/h} -- PUS-A and PUS-C secondary header handling,
verification support, source and destination identifiers, timestamps, and a
service/sub-service model. The active PUS version is configured in
\texttt{ecss\_pkt\_cfg.h}.\\

\item \texttt{ecss\_pkt\_cfg.h} -- configuration and service type /
sub-service definitions (ST\_VER, ST\_HK, ST\_EVT, ST\_Mem, ST\_Time, and so
on). The current configuration tracks PUS-C as the baseline and defines a
maximum packet size of 4096 bytes.\\

\item \texttt{crc16.c/h} -- CRC16-CCITT-FALSE with a slice-by-8 lookup table
and a benchmark suite.\\

\item \texttt{bitman.c/h} -- bit-level field extraction used by the packet
library.\\

\end{itemize}

\noindent
The \texttt{pus\_debug\_print()} function uses this library to print decoded
CCSDS and PUS-C headers and payloads for both directions when enabled with
\texttt{-D} and \texttt{-P}. The CRC16 check on the data path is performed
with \texttt{crc16\_buf()} from the same library (see
Section~\ref{sec:puscrc}).

\section{Block Diagram}

Figure~\ref{fig:block} shows the connectivity design of the bridge. Ground
support equipment, a grmon debug monitor or a \gls{GRESB} device exchange data
with the bridge over Ethernet (\gls{TCP}/IP or \gls{UDP}). The network side of
the bridge frames the received byte streams according to the selected packet
mode, the relaying core hands the packets to the \gls{SpW} side, which submits
them through a \gls{STAR} API device onto the \gls{SpaceWire} link, and the
same path is used in reverse for packets received from the link.

\begin{figure}[htbp]
\centering
\resizebox{0.9\textwidth}{!}{%
\begin{tikzpicture}[
	box/.style = {draw, rectangle, rounded corners = 2pt, fill = uvie-gray!25,
		      text width = 2.2cm, align = center, inner sep = 3pt},
	frame/.style = {draw, rectangle, rounded corners = 2pt, dashed,
			text width = 9.6cm, minimum height = 1.8cm, align = center}
]
	\node[frame] (bridge) at (8.2, 0) {};

	\node[box] (peers)  at (0, 0) {ground support equipment,\\ grmon};
	\node[box] (net)    at (4.4, 0) {net side\\ TCP/UDP\\ packet framing};
	\node[box] (core)   at (8.2, 0) {packet\\ relaying};
	\node[box] (spw)    at (12.0, 0) {SpW side\\ STAR API\\ device};
	\node[box] (target) at (16.2, 0) {SpW target\\ LEON3-FT,\\ GR712RC};

	\draw[->] (peers.east) -- (net.west);
	\draw[->] (net.east) -- (core.west);
	\draw[->] (core.east) -- (spw.west);
	\draw[->] (spw.east) -- (target.west);

	\node at (2.2, 0.45) {\scriptsize Ethernet};
	\node at (14.1, 0.45) {\scriptsize SpaceWire};

	\node at (8.2, -1.45) {\texttt{spw\_bridge}};
\end{tikzpicture}}
\caption{Block diagram of the \texttt{spw\_bridge} connectivity}
\label{fig:block}
\end{figure}

\section{Data Flow}

In the \emph{network to SpaceWire} direction, a poll thread waits for data on
the data service, frames it according to the selected mode, logs it, and hands
it to \texttt{spw\_send\_packet} which transmits it on the \gls{SpW} link (see
Figure~\ref{fig:block}). The
rmap service instead parses a pseudo-\gls{RMAP} command header (see
Section~\ref{sec:rmapif}) and issues the corresponding read or write command.
\\

In the \emph{SpaceWire to network} direction, the \gls{SpW} polling thread
receives a packet and calls \texttt{net\_pkt\_sink}. Depending on the
configuration, the packet is logged, its PUS sequence is checked, and any
RMAP replies (identified by the RMAP protocol identifier byte after the
skipped header) are forwarded to the RMAP service clients. The data is then
forwarded to the data service clients, optionally wrapped in a GRESB packet.


\chapter{Using the Bridge}
\label{cha:usage}

\section{Synopsis}

\begin{lstlisting}[language=bash]
  $ ./spw_bridge [OPTIONS]
\end{lstlisting}

\noindent
The bridge runs until it is interrupted; \texttt{SIGINT}, \texttt{SIGTERM}
and \texttt{SIGHUP} terminate the program and shut down the device cleanly.
The full option set is:

\begin{tabu} to \linewidth {@{}ll@{}}
	\texttt{-i DEVNUM}	& \gls{SpW} device id to use (default 0) \\
	\texttt{-c CHANNEL}	& \gls{SpW} channel to use (default 1) \\
	\texttt{-n PATH}	& routing path, colon-separated hex bytes, one byte per node \\
	\texttt{-p PORT}	& local data port number (default 1234) \\
	\texttt{-s ADDRESS}	& local source address (default 0.0.0.0) \\
	\texttt{-r HOST:PORT}& client mode: address and port of the remote target \\
	\texttt{-u}		& use UDP instead of TCP \\
	\texttt{-d NUM}		& number of header bytes to drop from incoming \gls{SpW} packets (default 0) \\
	\texttt{-t USEC}	& throttle: delay in microseconds inserted before each \gls{SpW} packet (default 0) \\
	\texttt{-S MBPS}	& link speed in Mbit/s (default 10, valid range 2--400) \\
	\texttt{-L LINKID}	& id of the link whose speed is set (default: equal to the channel) \\
	\texttt{-P}		& parse the network byte stream for \gls{PUS} packets \\
	\texttt{-D}		& print decoded PUS-C headers and payloads (requires \texttt{-P}) \\
	\texttt{-N}		& suppress payload bytes in the debug printout (short form) \\
	\texttt{-C}		& disable the PUS CRC16 check (see Section~\ref{sec:puscrc}) \\
	\texttt{-F}		& parse the network byte stream for FEE data packets \\
	\texttt{-R [PORT]}	& enable the pseudo--\gls{RMAP} service on PORT (default 2345) \\
	\texttt{-M C1:C2}	& monitor mode: bridge the given \gls{SpW} channels, copying packets verbatim between them \\
	\texttt{-E}		& decode \gls{RMAP} packets in the debug printout (requires \texttt{-M} and \texttt{-D}) \\
	\texttt{-G}		& use the GRESB protocol for the network exchange \\
	\texttt{-X}		& perform a device reset before opening the channel \\
	\texttt{-h}, \texttt{--help}& print the usage message and exit \\
\end{tabu}

\noindent
A routing path given with \texttt{-n} is a colon-separated list of hex byte
values, one per node (e.g. \texttt{-n 01:24:00:00:00}); up to 256 nodes are
accepted and out-of-range values are truncated to 8 bits.

\section{Operating Modes}

The bridge supports three modes of operation:

\begin{itemize}

\item \emph{server mode} (default): the bridge listens on the configured
address and port for an arbitrary number of TCP connections. Every connected
peer receives the \gls{SpW} traffic, and data from any peer is transmitted on
the \gls{SpW} link.\\

\item \emph{client mode} (\texttt{-r}): the bridge connects to the specified
\texttt{HOST:PORT} and operates a single connection. If the link is lost,
the bridge terminates.\\

\item \emph{UDP mode} (\texttt{-u}): the bridge binds a UDP socket and sends
one datagram per \gls{SpW} packet. Peers are registered as they are seen,
i.e. the first datagram from a host registers it as a recipient.\\

\end{itemize}

\section{SpaceWire Configuration}

At startup the bridge lists the detected \gls{STAR} devices and their names.
The device is selected with \texttt{-i} (default 0). The channel is selected
with \texttt{-c}; a non-existent channel is rejected and the available
channels are printed. The link signalling rate is set with \texttt{-S} to a
value between 2 and 400\,Mbit/s; the actual, measured receive rate is
reported during link start. \texttt{-L} selects the link whose speed is set;
on the Brick MK~II the transmit signalling rate of the second port is
configured in this way. \texttt{-X} resets the device before the channel is
opened; note that this affects other ports of the device as well.

\section{Interpreting the Network Stream}

By default the network byte stream is forwarded without framing (raw mode);
please note the packet merging and splitting caveat described in
Section~\ref{sec:net}. For reliable transport of packet-oriented traffic,
use one of the packet modes:

\begin{itemize}

\item \texttt{-P} frames the stream as \gls{PUS} packets, checks the CCSDS
sequence counter in both directions and verifies the PUS CRC16 (see
Section~\ref{sec:puscrc}). Sequence mismatches are reported to
{\ttfamily stderr}, together with the expected and received counter
values.\\

\item \texttt{-F} frames the stream as FEE data packets (protocol identifier
\texttt{0xF0}). All other packets fall back to raw framing.\\

\item \texttt{-G} frames the stream as GRESB protocol packets (see
Section~\ref{sec:gresb}).\\

\end{itemize}

\noindent
\texttt{-D} additionally prints a decoded CCSDS and PUS-C header and payload
dump of every packet passing through the bridge, in both directions. It
requires \texttt{-P}; once enabled, the printout can also be switched on and
off from the keyboard while the bridge is running (see
Section~\ref{sec:kbd}).
\\
\texttt{-N} restricts the printout to the decoded header lines and suppresses
the raw payload dump. Like the printout itself, the short form can be toggled
from the keyboard while the bridge is running (see Section~\ref{sec:kbd}).

\section{Monitor Mode}
\label{sec:monitor}

Monitor mode turns the bridge into a passive tap between two \gls{SpW}
channels of the same device. It is enabled with
\texttt{-M CHANNEL1:CHANNEL2}; the two colon-separated channel numbers select
the physical \gls{SpW} ports to be bridged. The two channels must differ.
\\
In this mode every packet received on one channel is copied verbatim to the
other, without adding or stripping a routing header and without any
interpretation on the forwarding path. The bridge neither inserts routing
bytes nor drops packet headers, and the forwarded \gls{SpW} transactions
remain atomic. The network side works as before, but it only observes: all
data received from the network is discarded, and the connected clients
receive a mirror of the copied \gls{SpW} traffic in both directions. The
channel selection \texttt{-c} and the header-drop option \texttt{-d} are
ignored; the routing path given with \texttt{-n} is not applied to the copied
packets. The link signalling rate set with \texttt{-S} is applied to both
channels.
\\
With \texttt{-D}, the debug printout shows the direction of the copied packet
as \texttt{SPW[A]->SPW[B]} instead of the network directions, followed by a
hex dump of the raw packet payload. With the additional \texttt{-E} option the
packets are decoded as \gls{RMAP} and a header line is printed in front of the
payload dump (see Section~\ref{sec:mondbg}). A packet that cannot be parsed as
\gls{RMAP} is reported as such. \texttt{-N} suppresses the payload dump (the
short form); the direction line and, with \texttt{-E}, the decoded
\gls{RMAP} header line are still printed. No \gls{PUS} interpretation is
performed on the monitor path, and packets are never modified by the
decoding.

\section{Reading the Monitor Mode Printout}
\label{sec:mondbg}

In monitor mode, the debug printout of every copied packet is preceded by a
time stamp and the direction of travel, using the actual channel numbers of the
bridge (e.g. \texttt{SPW[1]->SPW[2]} for a packet received on channel~1 and
copied to channel~2). Without \texttt{-E} the packet is shown as a raw hex dump
of its payload:

\begin{lstlisting}
[2026-09-10 10:20:30.512] SPW[1]->SPW[2]
  payload (18 bytes):
  00 01 02 03 04 05 06 07 08 09 0a 0b 0c 0d 0e 0f
  10 11
\end{lstlisting}

\noindent
With \texttt{-E}, a decoded \gls{RMAP} header line is printed in front of the
payload dump; a write command and the matching read reply, for example, look
like

\begin{lstlisting}
[2026-09-10 10:20:30.512] SPW[1]->SPW[2]
  RMAP: CMD WRITE dst=0xfe key=0xab src=0x30 tr_id=1 addr=0x00001000 data_len=2 hdr_crc=0x1b data_crc=0x4a
  payload (19 bytes):
  00 01 02 03 04 05 06 07 08 09 0a 0b 0c 0d 0e 0f
  10 11 12

[2026-09-10 10:20:32.013] SPW[2]->SPW[1]
  RMAP: REPLY READ dst=0x30 key=0x00 src=0xfe tr_id=1 addr=0x00001000 data_len=5 hdr_crc=0x77 data_crc=0x2c
  payload (22 bytes):
  00 01 02 03 04 05 06 07 08 09 0a 0b 0c 0d 0e 0f
  10 11 12 13 14 15
\end{lstlisting}

\noindent
The printed fields are

\begin{tabu} to \linewidth {@{}ll@{}}
	\texttt{SPW[A]->SPW[B]}	& direction of travel: the channel the packet was received on, and the channel it is copied to \\
	\texttt{CMD}, \texttt{REPLY} & whether the packet is an \gls{RMAP} command or a reply \\
	\texttt{WRITE}, \texttt{READ} & the operation of an \gls{RMAP} command \\
	\texttt{dst}, \texttt{src}	& target and initiator logical addresses \\
	\texttt{key}		& command authorisation key (status field in replies) \\
	\texttt{tr\_id}		& transaction identifier \\
	\texttt{addr}		& (first) data address \\
	\texttt{data\_len}	& length of the \gls{RMAP} data field in octets \\
	\texttt{hdr\_crc}, \texttt{data\_crc} & the header and data CRC-8 values carried by the packet \\
\end{tabu}

\noindent
A packet that cannot be parsed as \gls{RMAP} is announced with a
{\ttfamily not an RMAP packet} line and only its payload is dumped. In the
short form (\texttt{-N}) the payload dump is suppressed: only the time stamp,
the direction and, with \texttt{-E}, the decoded \gls{RMAP} header line are
printed, and a packet that is not \gls{RMAP} is not announced at all. The
payload itself is formatted as in the PUS printout (Section~\ref{sec:pusdbg}):
hex octets, sixteen per line, with an empty line closing the dump.

\section{Keyboard Control of the Debug Printout}
\label{sec:kbd}

When \texttt{-D} is active, the decoded debug printout can be toggled from the
terminal while the bridge is running: pressing \texttt{d} or \texttt{D} turns
it on or off, and each change is announced on the terminal with a
{\ttfamily PUS debug enabled} or {\ttfamily PUS debug disabled} message. This
allows a quick look at the PUS traffic without restarting the bridge or
resubmitting the command line. Likewise, \texttt{s} or \texttt{S} toggles the
short form of the printout, which suppresses the payload dump; the change is
announced with a {\ttfamily short debug form enabled} or
{\ttfamily short debug form disabled} message. The keyboard control is handled
by a dedicated monitor thread and requires \texttt{stdin} to be connected to a
terminal; if the bridge runs in the background or with redirected input, the
toggle is not available.

\section{Reading the PUS Debug Printout}
\label{sec:pusdbg}

With \texttt{-D} and \texttt{-P}, every packet is preceded by a time stamp and
the direction of travel, followed by a decoded CCSDS and PUS-C header dump and
a hex dump of the payload. An example for a telecommand (TC) with the debug
output enabled is

\begin{lstlisting}
[2026-09-10 10:20:30.512] NET->SPW
  CCSDS: type=TC APID=0x101 seq_flags=3 seq_count=42 data_len=35
  PUS-C: ST=5 SST=1 source_id=0x0000 ACK=0b0110
  CRC=OK
  payload (30 bytes):
  00 01 02 03 04 05 06 07 08 09 0a 0b 0c 0d 0e 0f
  10 11 12 13 14 15 16 17 18 19
\end{lstlisting}

\noindent
and for the corresponding telemetry (TM) return traffic

\begin{lstlisting}
[2026-09-10 10:20:31.048] SPW->NET
  CCSDS: type=TM APID=0x088 seq_flags=3 seq_count=7 data_len=31
  PUS-C: ST=1 SST=33 MTC=4 DEST_ID=0x0005 T_REF_STS=0010 TS=5c:00:12:34:00:42
  CRC=OK
  payload (18 bytes):
  00 01 02 03 04 05 06 07 08 09 0a 0b 0c 0d 0e 0f
  10 11
\end{lstlisting}

\noindent
The printed fields are

\begin{tabu} to \linewidth {@{}ll@{}}
	\texttt{APID}		& application process identifier of the CCSDS primary header \\
	\texttt{type}		& \texttt{TM} or \texttt{TC} \\
	\texttt{seq\_flags}, \texttt{seq\_count} & CCSDS sequence flags and sequence count (decimal) \\
	\texttt{data\_len}	& CCSDS data length in octets \\
	\texttt{ST}, \texttt{SST}& PUS service type and subservice type \\
	\texttt{source\_id}	& TC source identifier \\
	\texttt{ACK}		& TC acknowledgement flags, printed as four bits \\
	\texttt{MTC}		& TM message type counter \\
	\texttt{DEST\_ID}	& TM destination identifier \\
	\texttt{T\_REF\_STS}	& TM time reference status, printed as four bits \\
	\texttt{TS}		& TM time, six-octet PUS-C time of day \\
	\texttt{CRC}		& result of the CRC16 check (\texttt{OK} or \texttt{NOK}, see Section~\ref{sec:puscrc}) \\
\end{tabu}

\noindent
The payload is dumped as hex octets, sixteen per line, with an empty line
closing the dump. A payload that ends with a partial line is terminated
properly, so every packet dump is followed by a blank line.
\\
In monitor mode (Section~\ref{sec:monitor}) the dump shows the raw copied
packet; with \texttt{-E} an \gls{RMAP} header line (command or reply, read or
write, addresses, transaction identifier, data length and the CRCs) is printed
in its place.

\section{The PUS CRC16 Check}
\label{sec:puscrc}

Every \gls{PUS} packet carries a CRC16 field, as mandated by
\gls{ECSS}-E-ST-70-41C. With \texttt{-P}, the bridge verifies this field on
every packet in both directions and reports a failed check on
{\ttfamily stderr}. The result is also printed in the debug dump as a
\texttt{CRC=OK} or \texttt{CRC=NOK} line directly before the payload. The
bridge never drops a packet with a failed check, it only reports it; it is a
mediator on the data path. \texttt{-C} disables the check, which also
suppresses the corresponding printout line. Disable it only when a
mission-specific tailoring of the stream omits the CRC16 field.

\section{The Pseudo-RMAP Interface}
\label{sec:rmapif}

With \texttt{-R} (optionally followed by a port number, default 2345) the
bridge starts a second TCP server that implements a simplified \gls{RMAP}
command interface. Each command has the form

\begin{lstlisting}[language=bash]
  <spw addr><op>:<addr>:<size>:<data>
\end{lstlisting}

\noindent
encoded as a ten-byte header, followed by the write payload for write
operations:

\begin{tabu} to \linewidth {@{}ll@{}}
	byte 0		& target \gls{SpW} logical address \\
	byte 1		& operation: 0 = read, non-zero = write \\
	bytes 2--5	& 32-bit starting address, big-endian \\
	bytes 6--9	& data size, big-endian \\
	bytes 10...	& write payload (write operations only) \\
\end{tabu}

\noindent
A read command issues a read of 5 bytes at the given address; the size field
is currently ignored for reads. A write command writes the given payload with
address increment. The RMAP commands are assembled with the STAR RMAP packet
library, using a source logical address of \texttt{0x30} and the command key
\texttt{0xab} without an extended address. The RMAP replies produced by a
target (returned to \texttt{0x30}) are detected on the \gls{SpW} link by their
RMAP protocol identifier and forwarded to the connected RMAP clients in the
same ten-byte-plus-data format.
\\

For example, using netcat, a read of 5 bytes at address \texttt{0x1000} of a
target with logical address \texttt{0xfe} is issued with

\begin{lstlisting}[language=bash]
  $ printf '\xfe\x00\x00\x00\x10\x00\x00\x00\x00\x00' |
      nc localhost 2345
\end{lstlisting}

\noindent
and a two-byte write of \texttt{0xde 0xad} to the same address with

\begin{lstlisting}[language=bash]
  $ printf '\xfe\x01\x00\x00\x10\x00\x00\x00\x00\x02\xde\xad' |
      nc localhost 2345
\end{lstlisting}

\section{Examples}

\subsection{Server Mode with a PCIe Card}

Transmit and receive on channel 2 of the first detected device:

\begin{lstlisting}[language=bash]
  $ ./spw_bridge -c 2
\end{lstlisting}

\subsection{Link Monitoring}

To observe the traffic between two \gls{SpW} targets connected to channels 1
and 2 of the same device, copy the packets verbatim between the channels and
decode them as \gls{RMAP}:

\begin{lstlisting}[language=bash]
  $ ./spw_bridge -M 1:2 -D -E
\end{lstlisting}

\noindent
The network clients see a mirror of both directions of the copied traffic.

\subsection{Brick MK II Configuration}

The Brick MK~II channel is always 1. Configure the signalling rate of link 2
and send packets via the route with id 2 to the remote \gls{SpW} device with
node address \texttt{0x14}:

\begin{lstlisting}[language=bash]
  $ ./spw_bridge -c 1 -n 2:14 -L 2
\end{lstlisting}

\noindent
The same configuration in client mode, connecting to
\texttt{localhost:1234}:

\begin{lstlisting}[language=bash]
  $ ./spw_bridge -c 1 -n 2:14 -L 2 -r localhost:1234
\end{lstlisting}

\subsection{Use in the CHEOPS EGSE}

With the real data processing unit (DPU) connected to the \gls{EGSE} Brick,
use link 1 of the Brick as follows:

\begin{lstlisting}[language=bash]
  $ ./spw_bridge -c 1 -n 01:24:00:00:00 -d 4 -L 1 -p 5573 -P
\end{lstlisting}

\noindent
and link 2 of the Brick accordingly:

\begin{lstlisting}[language=bash]
  $ ./spw_bridge -c 1 -n 02:24:00:00:00 -d 4 -L 2 -p 5573 -P
\end{lstlisting}

\noindent
The first byte of the route is consumed by the Brick as the link number, the
next byte is the address of the DPU (\texttt{0x24}) and the remaining three
bytes are consumed by the DPU. The \texttt{-d 4} option drops the four-byte
\gls{SpW} address header that the real DPU attaches to its packets. Traffic is
sent to local port 5573 in \gls{PUS} framing mode.

\subsection{Use in the ARIEL Mission}

For the ARIEL mission, using a PCIe MK~II card, the bridge is typically
started as follows:

\begin{lstlisting}[language=bash]
  $ ./spw_bridge -i0 -P -n52:02:00:00 -p5573 -S60 -d4 -c3
\end{lstlisting}

\noindent
This selects device number 0 (\texttt{-i 0}), parses the network byte stream
as \gls{PUS} packets (\texttt{-P}), uses the routing path \texttt{52:02:00:00}
(\texttt{-n}), listens on local port 5573 (\texttt{-p 5573}), sets the link
signalling rate to 60\,Mbit/s (\texttt{-S 60}), drops the four-byte header of
incoming \gls{SpW} packets (\texttt{-d 4}) and opens \gls{SpW} channel 3
(\texttt{-c 3}). The channel flag selects the physical \gls{SpW} port of the
card and may need adjustment depending on which port is wired to the target.

\subsection{GRESB Protocol with grmon}
\label{sec:gresb}

For \gls{RMAP} access through the grmon debug monitor, connect the target
\gls{SpW} link to a \gls{SpW} port of the used device that supports \gls{RMAP}
(on the \gls{GR712RC} only link 0 or link 1 support \gls{RMAP}). Start the
bridge with the GRESB option and set the link speed to 10\,Mbit/s:

\begin{lstlisting}[language=bash]
  $ ./spw_bridge -c 1 -S 10 -G
\end{lstlisting}

\noindent
The GRESB uses a base port number of 3000 with two ports per virtual link
(the even port for transmit, the odd port for receive). Verify that these
ports are forwarded to the data port of the bridge, e.g. with socat:

\begin{lstlisting}[language=bash]
  $ socat tcp-l:3000,fork,reuseaddr tcp:127.0.0.1:1234
  $ socat tcp-l:3001,fork,reuseaddr tcp:127.0.0.1:1234
\end{lstlisting}

\noindent
You may now start grmon against the machine that runs the bridge:

\begin{lstlisting}[language=bash]
  $ grmon -gresb 127.0.0.1
\end{lstlisting}

\noindent
There is a caveat: the GRSPW2 \gls{SpW} port of the target may not be
initialised if no boot ROM is present. For the \gls{GR712RC} you have to
connect through the JTAG or another debug interface first and initialise the
core, e.g. GRSPW0:

\begin{lstlisting}[language=bash]
  grmon> wmem 0x80100800 0xA0010006
  grmon> wmem 0x8010080C 0x00000909
\end{lstlisting}

\noindent
For details refer to the GRMON2 manual \cite{GRMON2UM}.

\end{document}